盈利预测的核心不是把H1预告放大，而是把可验证的交付、价格/结构、确认、毛利、费用、税和少数股东归属连接起来。冻结基础情景预计2026--2028收入1,775/2,000/2,150亿元，归母172.76/216.49/246.81亿元，EPS 2.296/2.877/3.280元。2026接近Street，长期显著更谨慎。

\section{分析拆分：Growth与Base只用于驱动，不赋独立倍数}

2025A Growth业务定义为船舶造修及海工，收入1,312.79亿元、毛利153.86亿元、毛利率11.72\%。Base业务206.99亿元明确是\textbf{合并营业收入减Growth后的分析余项}：公司披露的配套、机电及其他主营收入186.17亿元，加未列入主营分产品表的其他营业收入20.82亿元。两者合计与合并利润表勾稽，但Base不是公司披露的新分部，也不能被表述为独立“配套分部”。

\begin{exhibitbox}[2025A分析基线]
\small
\begin{tabularx}{\exhibitboxwidth}{L{3.4cm}R{2.6cm}R{2.6cm}R{2.3cm}X}
\toprule
\textbf{分析业务} & \textbf{收入} & \textbf{毛利} & \textbf{毛利率} & \textbf{用途与边界} \\
\midrule
Growth：船舶造修及海工 & 1,312.79 & 153.86 & 11.72\% & 由交付、价格/结构和确认驱动；不拆独立EPS \\
Base：合并收入减Growth的分析余项 & 206.99 & 37.34 & 18.04\% & 186.17主营配套/机电/其他+20.82其他营业收入；不冒充披露分部 \\
合并 & 1,519.78 & 191.19 & 12.58\% & 与合并营业收入和成本勾稽 \\
\bottomrule
\end{tabularx}
\sourcenote{CF-03；206.99=186.17+20.82亿元。Base/Growth拆分为分析口径，不是公司披露的新分部。}
\end{exhibitbox}

这个拆分的“所以呢”是把订单驱动与相对稳定业务分开，同时避免PS/PEG滥用。费用、税、少数股东与投资收益仍在合并层面建模；没有Growth standalone EPS，也没有Growth独立目标价。

\section{E11｜从分业务毛利到归母净利润}

\begin{equation*}
\begin{split}
\text{合并毛利}_{t}&=\text{Growth收入}_{t}\times\text{Growth毛利率}_{t}\\
&\quad+\text{Base收入}_{t}\times\text{Base毛利率}_{t},\\
\text{归母净利}_{t}&=[\text{毛利}-\text{经营费用}-\text{减值}+\text{财务/投资/其他净收益}]\\
&\quad\times(1-\text{税率})\times(1-\text{少数股东份额}).
\end{split}
\end{equation*}

2026 Base Growth收入1,552.41亿元、毛利率16.8\%，Base收入222.59亿元、毛利率16.0\%，合并毛利296.42亿元。再扣经营费用128亿元、减值3亿元，加财务/投资/其他净收益36.5亿元并考虑税与归属后，归母172.76亿元。所有中间项为冻结研究假设，必须随正式中报整体重跑。

\begin{exhibitbox}[Base分部到合并盈利桥]
\small
\begin{tabularx}{\exhibitboxwidth}{L{3.6cm}R{2.4cm}R{2.4cm}R{2.4cm}X}
\toprule
\textbf{项目} & \textbf{2026E} & \textbf{2027E} & \textbf{2028E} & \textbf{验证重点} \\
\midrule
Growth收入 & 1,552.41 & 1,779.28 & 1,916.74 & 交付、价格/结构、确认 \\
Base收入 & 222.59 & 220.72 & 233.26 & 配套与其他业务稳定性 \\
合并收入 & 1,775.00 & 2,000.00 & 2,150.00 & 正式收入及分产品披露 \\
合并毛利率 & 16.70\% & 17.30\% & 17.52\% & 高价订单持续期 \\
经营费用 & 128.00 & 137.00 & 136.00 & 整合效率与费用口径 \\
减值 & 3.00 & 3.00 & 4.00 & 返工、合同与资产质量 \\
财务/投资/其他净收益 & 36.50 & 47.00 & 55.00 & 不把单季收益永久化 \\
税前利润（PBT） & 201.92 & 253.03 & 291.64 & 毛利减费用/减值并加净收益 \\
模型税率/所得税 & 8.0\% / 16.15 & 8.0\% / 20.24 & 9.0\% / 26.25 & 冻结情景税率，不是公司指引 \\
合并净利润 & 185.77 & 232.78 & 265.39 & PBT减所得税 \\
少数股东份额/损益 & 7.0\% / 13.00 & 7.0\% / 16.29 & 7.0\% / 18.58 & 由合并净利到归母 \\
归母净利润 & 172.76 & 216.49 & 246.81 & 税、投资收益与少数股东 \\
EPS & 2.296 & 2.877 & 3.280 & 75.2562亿股分母 \\
\bottomrule
\end{tabularx}
\sourcenote{冻结增长模型；金额单位除EPS与比例外均为亿元。所得税=PBT×税率；少数股东损益=合并净利润×少数股东份额。}
\end{exhibitbox}

Base完整桥说明归母并不是“毛利减费用”的快捷结果：2026--2028模型税率为8\%/8\%/9\%，少数股东份额均为7\%。Bear税率为10\%/12\%/13\%、少数股东份额为15\%/18\%/20\%；Bull税率为7\%/8\%/8\%、少数股东份额均为4\%，完整数值见模型附录。估值风险仍集中在毛利持续期与归属：2026 Base Growth毛利率16.8\%低于Q1合并毛利率，避免把单季高点永久化，却仍较2025核心业务11.72\%明显上升；2027--2028继续改善到17.5\%/17.8\%，是最重要的可证伪假设。

\section{E14｜Bear/Base/Bull完整经营情景}

Bear不是需求立即消失，而是交付延误、价格/结构回落、设备与返工成本侵蚀以及现金占用共同发生；Bull则要求交付超计划、结构与效率同时改善，并且现金回收跟上。三情景不包含军品、沪东中华、集团未分配订单或弱来源目标价。

\begin{exhibitbox}[2026--2028三情景结果]
\scriptsize
\begin{tabularx}{\exhibitboxwidth}{L{1.8cm}R{2.0cm}C{1.8cm}R{2.0cm}R{1.8cm}R{1.8cm}Y}
\toprule
\textbf{情景/年} & \textbf{收入} & \textbf{综合GM} & \textbf{归母} & \textbf{EPS} & \textbf{现价PE} & \textbf{OCF/NP} \\
\midrule
Bear 2026 & 1,652.49 & 14.19\% & 112.07 & 1.489 & 22.17x & 0.55x \\
Bear 2027 & 1,641.28 & 13.76\% & 93.72 & 1.245 & 26.52x & 0.65x \\
Bear 2028 & 1,666.22 & 13.10\% & 80.20 & 1.066 & 30.99x & 0.75x \\
\rowcolor{accentblue!8} Base 2026 & 1,775.00 & 16.70\% & 172.76 & 2.296 & 14.38x & 0.85x \\
\rowcolor{accentblue!8} Base 2027 & 2,000.00 & 17.30\% & 216.49 & 2.877 & 11.48x & 0.95x \\
\rowcolor{accentblue!8} Base 2028 & 2,150.00 & 17.52\% & 246.81 & 3.280 & 10.07x & 1.00x \\
Bull 2026 & 1,935.00 & 18.27\% & 217.34 & 2.888 & 11.43x & 1.05x \\
Bull 2027 & 2,205.00 & 19.22\% & 279.85 & 3.719 & 8.88x & 1.10x \\
Bull 2028 & 2,448.00 & 19.91\% & 331.63 & 4.407 & 7.49x & 1.15x \\
\bottomrule
\end{tabularx}
\sourcenote{冻结增长模型。金额单位：亿元；EPS为元/股。}
\end{exhibitbox}

情景表对行动的含义是明显的非对称性：Base已经接近现价，Bear在盈利和倍数上同时受压。现价不是低价买入周期复苏的简单赔率，而是要求模型至少达到Base并逐步证明现金质量。

\section{H1残余与二维敏感性}

若正式H1归母分别为92/101/110亿元，2026 Base全年172.76亿元要求H2实现80.76/71.76/62.76亿元。该残余分析承认交付时点性，不采用H1乘二。Q1与H1只是阶段锚，完整全年仍由交付和毛利桥决定。

\begin{exhibitbox}[2026归母对确认系数与Growth毛利率敏感性]
\small
\begin{tabularx}{\exhibitboxwidth}{L{3.6cm}R{3.0cm}R{3.0cm}R{3.0cm}X}
\toprule
\textbf{确认系数\textbackslash Growth GM} & \textbf{15.8\%} & \textbf{16.8\%} & \textbf{17.8\%} & \textbf{解释} \\
\midrule
0.92 & 150.74 & 163.46 & 176.19 & 确认偏慢 \\
0.96 & 159.48 & \textbf{172.76} & 186.04 & Base确认 \\
1.00 & 168.22 & 182.06 & 195.90 & 确认偏快 \\
\bottomrule
\end{tabularx}
\sourcenote{冻结增长模型；其余变量保持2026 Base。金额单位：亿元。}
\end{exhibitbox}

敏感性说明1个百分点毛利变化与确认节奏都足以改变归母十亿元以上。正式半年报必须同时观察收入、核心毛利、合同资产和现金；只看归母总额无法区分毛利改善与确认前置。

\section{模型许可与零信用清单}

当前模型状态为CONDITIONAL，但可复算性为PASS。允许的信用是合并口径年度盈利；禁止船型/军品/船厂SOTP、Growth独立高增长倍数、在手直接折现和子公司利润重复加总。中远项目只做去重订单质量证据。

\begin{keyinsight}[盈利预测的“所以呢”]
2026利润方向可信，真正的预期差在2027--2028毛利持续期与现金。新增仓位只有在正式数据使Base盈利上修、同时不降低现金质量时才合理；单独上修订单、交付或H1利润，不足以改变33.07元目标。
\end{keyinsight}

预测风险因此可被明确监测：交付影响收入量，价格/结构与成本影响毛利，确认影响时点，少数股东与投资收益影响归属，OCF约束倍数。任何关键项脱离披露，都只能落入情景而不能写成事实。
